\chapter{Conclusion g\'en\'erale}

\section{Synth\`ese des travaux}

Ce m\'emoire a pr\'esent\'e la conception, le d\'eveloppement et la validation d'un syst\`eme Andon intelligent bas\'e sur l'Internet Industriel des Objets (\iiot) pour la gestion en temps r\'eel des pannes de production chez \sews.

Le travail r\'ealis\'e a mobilis\'e un large \'eventail de comp\'etences techniques : d\'eveloppement embarqu\'e (ESP32/Arduino), architecture r\'eseau (MQTT, Wi-Fi), d\'eveloppement backend (Node.js/Express), conception de base de donn\'ees (PostgreSQL) et d\'eveloppement frontend (Angular). L'ensemble des composants a \'et\'e int\'egr\'e de mani\`ere coh\'erente pour former un syst\`eme fonctionnel et performant.

\section{R\'ealisation des objectifs}

Les objectifs d\'efinis en d\'ebut de projet ont \'et\'e atteints :

\begin{enumerate}
    \item \textbf{Analyse du contexte :} l'immersion au sein de l'atelier Test \'Electrique de \sews a permis d'identifier pr\'ecis\'ement les besoins en mati\`ere de gestion des pannes.
    
    \item \textbf{Conception de l'architecture :} une architecture en quatre couches (perception, communication, traitement, pr\'esentation) a \'et\'e d\'efinie et valid\'ee.
    
    \item \textbf{Impl\'ementation :} chaque composant a \'et\'e d\'evelopp\'e : firmware ESP32, serveur Node.js, base de donn\'ees PostgreSQL et application Angular.
    
    \item \textbf{Validation :} les tests en conditions r\'eelles pendant 3 mois ont d\'emontr\'e les performances du syst\`eme.
\end{enumerate}

\section{Principaux r\'esultats}

Les r\'esultats obtenus confirment l'efficacit\'e du syst\`eme d\'evelopp\'e :

\begin{itemize}
    \item \textbf{Temps de d\'etection :} 189 ms en moyenne, inf\'erieur \`a la cible de 1 seconde.
    \item \textbf{Temps de r\'eponse :} r\'eduction de 33\% (de 18 \`a 12 minutes).
    \item \textbf{Taux de disponibilit\'e :} am\'elioration de 6,6 points (de 87,2\% \`a 93,8\%).
    \item \textbf{Fiabilit\'e :} 99,98\% de messages MQTT livr\'es sans perte.
    \item \textbf{Co\^ut :} mat\'eriel \`a 96EUR, infrastructure 696EUR/an.
\end{itemize}

\section{Difficult\'es rencontr\'ees}

La r\'ealisation de ce projet a confront\'e \`a plusieurs difficult\'es :

\begin{itemize}
    \item \textbf{Environnement industriel :} les perturbations \'electromagn\'etiques de l'atelier ont n\'ecessit\'e des pr\'e-cautions suppl\'ementaires pour les n\oe{}uds ESP32.
    \item \textbf{R\'eseau Wi-Fi :} la couverture et la stabilit\'e du r\'eseau sans fil dans un environnement industriel ont requis des optimisations.
    \item \textbf{Int\'e-gration :} la coordination entre les diff\'erents composants (mat\'eriel, backend, frontend) a n\'ecessit\'e une approche it\'erative.
    \item \textbf{Temps de d\'eveloppement :} la dur\'ee limit\'ee du stage (4 mois) a impos\'e des choix pragmatiques de prioritisation.
\end{itemize}

\section{Perspectives d'\'evolution}

Plusieurs pistes d'am\'elioration et d'\'evolution sont envisageables :

\subsection{Court terme}

\begin{itemize}
    \item \textbf{Mobile :} d\'eveloppement d'une application mobile native (React Native) pour un acc\`es mobile optimis\'e.
    \item \textbf{Offline :} ajout d'un mode d\'egrad\'e pour les n\oe{}uds ESP32 en cas de perte de connexion.
    \item \textbf{Multilingue :} internationalisation de l'interface en fran\c{c}ais, arabe et anglais.
\end{itemize}

\subsection{Moyen terme}

\begin{itemize}
    \item \textbf{Analyse pr\'edictive :} int\'e-gration de mod\`eles d'apprentissage automatique pour la pr\'ediction des pannes.
    \item \textbf{Capteurs suppl\'ementaires :} ajout de capteurs de temp\'erature, vibration et courant pour un diagnostic enrichi.
    \item \textbf{Extension multi-sites :} g\'en\'eralisation du syst\`eme \`a tous les sites de production SEWS.
\end{itemize}

\subsection{Long terme}

\begin{itemize}
    \item \textbf{Jumeau num\'erique :} cr\'eation d'une maquette num\'erique de l'atelier pour la simulation.
    \item \textbf{Int\'e-gration ERP :} connexion avec le syst\`eme d'information de l'entreprise (SAP, Oracle).
    \item \textbf{Microservices :} migration vers une architecture microservices pour la scalabilit\'e.
\end{itemize}

\section{Conclusion personnelle}

Ce stage a \'et\'e une exp\'erience enrichissante sur les plans technique et humain. Il m'a permis de mettre en pratique les connaissances acquises durant ma formation de Master, tout en d\'e-couvrant les r\'ealit\'es de l'environnement industriel. Le d\'eveloppement de ce syst\`eme Andon IIoT m'a offert l'opportunit\'e d'explorer des technologies modernes de l'IoT et de comprendre les enjeux de la transformation num\'erique dans l'industrie.

Je suis reconnaissant envers mon encadrant acad\'emique, Pr. Kaoutar ALLABOUCHE, et mon encadrant industriel, M. Mohamed EDDOUCH, pour leur accompagnement et leurs conseils tout au long de ce projet.
